% ══════════════════════════════════════════════════════════════════
% Section 5 — Fortran Export Pipeline
% ══════════════════════════════════════════════════════════════════
\section{Fortran export pipeline}
\label{sec:fortran_export}

A key contribution of this work is the automated export of trained neural
networks into self-contained Fortran~90 subroutines that can be linked
directly to commercial and research finite element solvers without any
Python or PyTorch dependency at runtime.

\subsection{Export strategy}
\label{sec:fortran_export:strategy}

We provide three export modes, targeting different use cases:

\begin{enumerate}
  \item \textbf{Standalone NN emitter} (\texttt{FortranEmitter}): Exports the
    neural network forward pass as a pure Fortran~90 subroutine.  Weights and
    biases are embedded as compile-time \texttt{PARAMETER} arrays.  This is
    useful for evaluating the energy outside a FE context.

  \item \textbf{Hybrid NN--UMAT emitter} (\texttt{HybridUMATEmitter}):
    Produces a complete UMAT subroutine where the neural network evaluates
    the strain energy function $\hat{\sef}$ and its first and second
    derivatives with respect to invariants, while all kinematic quantities
    (deformation tensors, invariant derivatives, push-forward to Cauchy
    stress, Jaumann correction) are computed analytically in Fortran.

  \item \textbf{Analytical UMAT emitter} (\texttt{UMATHandler}): Generates
    fully symbolic Fortran code from the material model using SymPy common
    subexpression elimination (CSE).  This serves as the reference
    implementation for validation.
\end{enumerate}

\subsection{Hybrid NN--UMAT architecture}
\label{sec:fortran_export:hybrid}

The hybrid UMAT follows the computational flow shown in
\Cref{fig:hybrid_umat}.  At each integration point, the solver provides the
deformation gradient $\Fmat$.  The UMAT:

\begin{enumerate}
  \item Computes $\Cmat$, $J$, and invariants $\{I_\alpha\}$ analytically.
  \item Normalizes invariants: $\hat{\tens{x}} = (\tens{x} - \bm{\mu}) / \bm{\sigma}$.
  \item Evaluates the NN forward pass to obtain $\hat{\sef}$.
  \item Backpropagates through the NN to obtain first derivatives
    $\partial\hat{\sef}/\partial\hat{x}_\alpha$ and second derivatives
    (Hessian) $\partial^2\hat{\sef}/\partial\hat{x}_\alpha\partial\hat{x}_\beta$.
  \item Transforms to physical gradients:
    $\partial\sef/\partial I_\alpha = (\partial\hat{\sef}/\partial\hat{x}_\alpha) / \sigma_\alpha$.
  \item Assembles the second Piola--Kirchhoff stress $\Smat$ via
    \Cref{eq:pk2_chain}.
  \item Pushes forward to Cauchy stress:
    $\sigmat = J^{-1}\Fmat\,\Smat\,\Fmat^T$.
  \item Assembles the spatial tangent $\mathbb{c}$ and applies the Jaumann
    correction for the Abaqus stress rate.
\end{enumerate}

The Hessian backpropagation through the NN layers is performed analytically in
Fortran using the chain rule through each layer's activation, avoiding any
need for an automatic differentiation library.

% ── Hybrid UMAT diagram ──────────────────────────────────────────
\begin{figure}[htbp]
\centering
% TODO: Replace with TikZ flowchart:
%   DFGRD1(F) -> C, J, invariants -> normalize -> NN forward -> W
%        -> NN backprop -> dW/dI, d2W/dIdJ
%        -> analytical S = 2 sum dW/dI * dI/dC
%        -> sigma = 1/J F S F^T
%        -> tangent + Jaumann -> DDSDDE
\fbox{\parbox{0.9\textwidth}{\centering\vspace{5cm}
  \textbf{[PLACEHOLDER]} \\[0.5em]
  Hybrid NN--UMAT computational flow: \\
  $\Fmat \to \Cmat, J, \{I_\alpha\}
   \to \text{normalize} \to \text{NN}(\hat{\sef})
   \to \text{backprop}(\partial\hat{\sef}/\partial I, \;
     \partial^2\hat{\sef}/\partial I^2)$ \\
  $\to \Smat = 2\sum \frac{\partial\sef}{\partial I_\alpha}
     \frac{\partial I_\alpha}{\partial\Cmat}
   \to \sigmat = J^{-1}\Fmat\Smat\Fmat^T
   \to \text{DDSDDE}$
\vspace{1cm}}}
\caption{Computational flow of the hybrid NN--UMAT subroutine.  Blue blocks
  are computed by the neural network; gray blocks are analytical Fortran.}
\label{fig:hybrid_umat}
\end{figure}

\subsection{Fortran implementation details}
\label{sec:fortran_export:details}

\paragraph{Weight embedding.}
All network weights and biases are written as Fortran \texttt{PARAMETER}
arrays with 15-digit double-precision formatting (\texttt{:.15e}).  This
avoids any file I/O at runtime and allows the compiler to optimize memory
layout.

\paragraph{Column-major storage.}
Weight matrices follow Fortran's column-major convention.  The forward pass
uses \texttt{MATMUL} intrinsics for matrix--vector products.

\paragraph{Activation functions.}
Each activation (softplus, tanh, sigmoid, ReLU) is implemented as an
element-wise Fortran loop.  For softplus, a numerically stable form
$\phi(x) = \log(1 + \exp(x))$ with overflow guards is used.

\paragraph{Supported architectures.}
The hybrid emitter supports MLP and PolyconvexICNN architectures.  For the
polyconvex model, each branch is emitted as a separate forward pass, and
the total energy is the sum of branch outputs.

\paragraph{Isotropic and anisotropic support.}
Isotropic UMATs use 3 invariant inputs ($\bar{I}_1, \bar{I}_2, J$);
anisotropic UMATs accept 5 inputs ($\bar{I}_1, \bar{I}_2, J, I_4, I_5$)
and include the fiber direction as a material constant.

% ── Code snippet ──────────────────────────────────────────────────
\begin{figure}[htbp]
\centering
% TODO: Replace with actual Fortran snippet from exported UMAT
%   showing the NN forward pass and stress assembly.
\fbox{\parbox{0.9\textwidth}{\centering\vspace{3cm}
  \textbf{[PLACEHOLDER]} \\[0.5em]
  Annotated Fortran~90 code excerpt from the hybrid NN--UMAT showing: \\
  (1) invariant computation, (2) NN forward pass with embedded weights, \\
  (3) backpropagation for $\partial\sef/\partial I$, (4) stress assembly.
\vspace{1cm}}}
\caption{Excerpt from an auto-generated hybrid NN--UMAT (Neo-Hookean, MLP
  $64 \times 64$).}
\label{fig:fortran_excerpt}
\end{figure}
